
\begin{table*}[t]
\caption{Closest methodological and domain precedents and the restricted comparison made in this study.}\label{tab:literature}
\centering\footnotesize\setlength{\tabcolsep}{3pt}\renewcommand{\arraystretch}{1.17}
\begin{tabularx}{\textwidth}{>{\hsize=0.78\hsize\linewidth=\hsize\raggedright\arraybackslash}X>{\hsize=1.0499999999999998\hsize\linewidth=\hsize\raggedright\arraybackslash}X>{\hsize=1.17\hsize\linewidth=\hsize\raggedright\arraybackslash}X}
\toprule
\textbf{Prior work} & \textbf{Relevant precedent} & \textbf{Question addressed here} \\ \midrule
Pinto et al.; Cheng et al.; Mishra et al. & Predictive timing features and feature/process comparisons & Magnitude of a fixed timing-feature addition under count CRPS within each family \\
SEISMIC; TiDeH; MaSEPTiDE & Event-history models of growth & Compact timing summaries with no marks or specialized process-model superiority claim \\
CASPER; Haimovich et al. & Conditional future moments and finite-horizon popularity & A fully specified integer-count distribution at a root-relative endpoint \\
DeepHawkes; CasFlow; CTCP; CasFT; ConCat & Temporal/structural representation learning & Interpretable restricted-input comparisons, evaluated using proper distribution scores \\
Peng et al. (2026 preprint) & Protocol-sensitive cascade benchmarking & Source-specific distributional information contrasts with explicit split limitations \\
QRF; NGBoost; count-regression literature & Established distribution families & Empirical behavior when the available prefix information changes \\
Martin et al.; Hofman et al. & Predictive accuracy and its limitations in social systems & Procedure-level gains, with no claim to identify intrinsic predictability \\
Count-assessment and proper-score literature & Calibration and distribution evaluation & Joint interpretation of CRPS, no-growth probabilities and interval behavior \\
\bottomrule\end{tabularx}
\end{table*}
